\documentclass[11pt]{scrartcl}

\usepackage[stixtwo]{fontsetup}
\usepackage{amsmath}

\begin{document}
\begin{enumerate}
	\item \textbf{Szöveges gondolatmenet}

	      A gyökjel alatt a számlálóban a leggyorsabban növekvő tag a
	      $3^{n + 2} = 9 \cdot 3^n$, hiszen $n^3 \ll 3^n$, a nevezőben pedig a
	      $7n$. Mivel $\sqrt[n]{c} \to 1$ és $\sqrt[n]{n} \to 1$, a határérték a
	      domináns exponenciális tag alapja, vagyis $3$ lesz.

	      \textbf{Becslés a rendőrelvvel}

	      $n \ge 1$ esetén $2n^3 \le 3^{n + 2}$, továbbá
	      $7n \le 2 + 7n \le 9n$. Becsüljük alulról és felülről a gyökjel
	      alatti törtet:
	      \[
		      \frac{3^{n + 2}}{9n}
		      \le \frac{3^{n + 2} + 2n^3}{2 + 7n}
		      \le \frac{2 \cdot 3^{n + 2}}{7n}
		      \text.
	      \]
	      Mindhárom kifejezésből $n$-edik gyököt vonva:
	      \[
		      3
		      \leftarrow \frac{3}{\sqrt[n]{n}}
		      \le \sqrt[n]{\frac{3^{n + 2} + 2n^3}{2 + 7n}}
		      \le \frac{3 \cdot \sqrt[n]{\frac{18}{7}}}{\sqrt[n]{n}}
		      \to 3
		      \text.
	      \]
	      A rendőrelv alapján tehát a határérték $3$.

	\item \textbf{Szöveges gondolatmenet}

	      A számlálóban a leggyorsabban növekvő tag az $n^n$, a nevezőben pedig
	      az $n!$, hiszen
	      \[
		      n^3 \ll 5^n \ll n! \ll n^n
		      \text,
	      \]
	      és $\cos^2(n)$ korlátos. Mivel a számláló domináns tagja
	      nagyságrendileg gyorsabban nő, a határérték $\infty$ lesz.

	      \textbf{Másik megoldás}

	      Osszuk el a számlálót és a nevezőt is $n!$-sal:
	      \[
		      \frac{5^n + n^n}{n^3 \cos^2(n) + n!}
		      = \frac{
			      \frac{5^n}{n!} + \frac{n^n}{n!}
		      }{
			      \frac{n^3 \cos^2(n)}{n!} + 1
		      }
		      \to \infty
		      \text,
	      \]
	      hiszen
	      \[
		      \frac{n^n}{n!} \ge n \to \infty,
		      \qquad
		      \frac{5^n}{n!} \to 0,
		      \qquad
		      0 \le \frac{n^3 \cos^2(n)}{n!} \le \frac{n^3}{n!} \to 0
		      \text.
	      \]
\end{enumerate}
\end{document}
